S'utilitza la lent d'un escàner per a capturar la imatge d'un full situat a 20,0~cm de la lent. La imatge que es forma en el sensor de l'escàner està invertida i té la mateixa mida que l'objecte.

\begin{apartat}{1,25}
A quina distància de la lent s'ha de col·locar el sensor perquè es formi una imatge nítida? Calculeu la distància focal imatge de la lent i indiqueu si es tracta d'una lent convergent o divergent.
\end{apartat}

\begin{apartat}{1,25}
Dibuixeu el diagrama de rajos corresponent en la quadrícula de sota indicant clarament totes les distàncies representades i utilitzades en l'apartat anterior.
\end{apartat}

\figura[0.85]{fig1.png}
